\section{Literature}
\label{sec:literature}

This paper sits at the intersection of four literatures: economic complexity and trade, trade costs and payment frictions, AML regulation and cross-border banking, and stablecoins and cross-border financial flows.

\paragraph{Economic complexity and trade.}
\citet{Hidalgo2007_product_space} establish that countries move through ``product space''---a network where proximity between products reflects shared capability requirements. Countries diversify into products close to their existing export basket, and this process shapes long-run development trajectories \citep{Hidalgo2009_building_blocks, Hausmann2014_atlas}. The Product Complexity Index (PCI) ranks products by the breadth and sophistication of capabilities they require: complex products (high PCI) are exported by few diversified economies, while simple products (low PCI) are ubiquitous \citep{Hidalgo2021_economic_complexity}. \citet{Hartmann2017_complexity_inequality} link economic complexity to income inequality, showing that structural transformation toward complex products generates more equitable growth. \citet{Bahar2014_neighbors} demonstrate that geographic proximity facilitates knowledge spillovers that enable complexity upgrading. \citet{Ndoya2024_financial_complexity} show that financial development supports complexity, but the specific role of payment infrastructure remains unexplored. None of this literature connects product complexity to the payment cost structure of international trade.

\paragraph{Trade costs and payment frictions.}
\citet{Anderson2004_trade_costs} provide a comprehensive taxonomy of trade costs, identifying transportation, border, and behind-the-border barriers. \citet{Head2014_gravity} survey the modern gravity toolkit, and \citet{SantosSilva2006_ppml} establish PPML as the consistent estimator for multiplicative trade models with heteroskedasticity and zeros. Within the broader trade-cost literature, a subset focuses on financial frictions. \citet{Manova2013_credit_constraints} shows that credit constraints reduce trade heterogeneously across firms sorted by financial vulnerability, and \citet{AntrasFolley2015_trade_finance} document how firms choose payment terms (cash-in-advance, open account, letters of credit) based on contracting frictions. \citet{Freund2004_internet_trade} and \citet{Fink2005_communication} show that communication technology reduces trade costs, providing a precedent for studying payment-technology effects.

\paragraph{AML regulation, correspondent banking, and harmonization.}
A separate literature studies how AML/CFT regulation reshapes the bilateral correspondent-banking network and, through it, cross-border payments. \citet{Rice2020_correspondent} document the 25\% decline in correspondent-banking relationships since 2012. \citet{Borchert2024_broken_relationships} provide the closest antecedent to our work: they show that de-risking causally reduces trade flows and identify FATF greylisting as a discrete trigger for correspondent-account closures. Their analysis treats the trade effect as uniform across products. We extend their framework in two directions. First, we interact the FATF greylist shock with product complexity, revealing that the trade damage concentrates in complex goods. Second, we add a friction-down counterpart---the EU AMLD harmonization shock---and show that complex trade rises in EU corridors after 2017 by a magnitude comparable to the loss from FATF greylisting. The symmetry across the two natural experiments isolates the payment-friction channel from confounders. AMLD-style harmonization is the policy mirror of greylisting: harmonization lowers the per-corridor compliance burden disproportionately for the products that pass through more compliance gates.

\paragraph{Stablecoins and cross-border flows.}
A growing literature examines cryptocurrency in the international financial architecture. \citet{GrafvonLuckner2023_capital_flows} show that crypto transactions follow economic fundamentals but bypass traditional intermediaries. \citet{GrafvonLuckner2024_capital_flight} document crypto as a channel for capital flight from countries with capital controls. \citet{Cerutti2024_crypto_flows} develop a measurement framework for cross-border crypto flows and identify their macro-financial drivers. \citet{Auer2025_defiying_gravity} estimate a gravity model of cross-border cryptocurrency flows, finding that crypto ``defies'' gravity: distance elasticities are lower than for traditional finance. \citet{FerrarMinesso2026_payment_links} provide the methodological bridge: they estimate the trade effects of interlinking national fast-payment systems, showing that payment integration increases bilateral goods trade. \citet{Rose2000_one_money} establishes the benchmark that common currencies triple trade---a result that frames the potential magnitude of stablecoin effects, since stablecoins effectively create a common settlement currency (USD) for participants regardless of national currency. \citet{Auer2021_multi_cbdc} discuss multi-CBDC arrangements as an alternative infrastructure for cross-border payments. \citet{IMF2025_stablecoins} provides the most recent IMF assessment of stablecoin macro-financial implications.

\paragraph{Our contribution.}
We bridge these four literatures. From the complexity literature, we take the insight that products differ in their capability requirements. From the trade-cost literature, we take the structural gravity framework. From the AML-regulation literature, we take FATF greylisting and EU AMLD as discrete identifying shocks. From the stablecoin literature, we take the premise that blockchain-based payment rails represent a viable alternative settlement layer. The resulting contribution is two-sided. First, two distinct natural experiments---one friction-up, one friction-down---identify a symmetric effect of payment frictions on complex trade. Second, the corridor-level estimates feed a product-level opportunity map identifying which country--product cells stand to gain most from reduced payment frictions and where deployment is feasible.
